% Generated draft: requires mathematical and source review.
\documentclass[11pt]{article}
\usepackage[margin=25mm]{geometry}
\usepackage{fontspec}
\setmainfont{Latin Modern Roman}
\usepackage{amsmath,amssymb,mathtools,unicode-math}
\setmathfont{Latin Modern Math}
\usepackage{enumitem,needspace,xurl,hyperref}
\setlength{\parindent}{0pt}
\setlength{\parskip}{5pt}
\setlength{\emergencystretch}{3em}
\newcommand{\N}{\mathbb{N}}
\newcommand{\Z}{\mathbb{Z}}
\newcommand{\Q}{\mathbb{Q}}
\newcommand{\R}{\mathbb{R}}
\newcommand{\C}{\mathbb{C}}
\newcommand{\F}{\mathbb{F}}
\newcommand{\A}{\mathbb{A}}
\newcommand{\PP}{\mathbb{P}}
\newcommand{\E}{\mathbb{E}}
\newcommand{\eps}{\varepsilon}
\newcommand{\dd}{\,\mathrm{d}}
\begin{document}
\section*{UnsolvedMath ID: 700335}
\textbf{Source problem code:} catalogue entry
\subsection{Definitions and mathematical statement}
For a natural-number degree \(n\), let \(\mathbb{A}_n\) denote the group of even permutations of a set of \(n\) points. A group word \(u(x,y,z)\) is evaluated at \(\xi,\eta,\zeta\) by substituting these elements for its variables. The source requires \(u\) to be a ``commutator word'' but does not specify the precise class of words meant. For example, membership in the derived subgroup of the free group on \(x,y,z\) and being an iterated commutator are different possible requirements; neither is imposed here without qualification.

The requested assertion is: for every triple of pairwise distinct primes \(p,q,r\) and every commutator word \(u(x,y,z)\) in the intended, unspecified sense, there exist a natural number \(n\) and elements \(\xi,\eta,\zeta\in\mathbb{A}_n\) such that
\[
\langle\xi,\eta,\zeta\rangle=\mathbb{A}_n,
\qquad
\xi^p=\eta^q=\zeta^r=1,
\qquad
\xi\eta\zeta\,u(\xi,\eta,\zeta)=1.
\]
Here \(1\) is the identity permutation. The power relations require orders dividing \(p,q,r\), respectively; they do not explicitly require exact orders or nonidentity generators.

The source also asks, parenthetically, whether for every fixed admissible choice of \(p,q,r,u\), the set
\[
\left\{n\in\N:\ \exists\xi,\eta,\zeta\in\mathbb{A}_n,\quad
\langle\xi,\eta,\zeta\rangle=\mathbb{A}_n,\quad
\xi^p=\eta^q=\zeta^r=\xi\eta\zeta\,u(\xi,\eta,\zeta)=1
\right\}
\]
is infinite. This is retained as a stronger question, rather than asserted as part of the basic existence claim.

No lower bound such as \(n\geq5\), exclusion of trivial alternating groups, or convention concerning whether \(0\in\N\) is supplied. These boundary conditions and the meaning of ``commutator word'' require clarification before the intended nontrivial formulation is fixed.
\subsection{Short English statement}
Given three pairwise distinct primes and a commutator word, ask whether some alternating group has three generators satisfying the specified power relations and the relation \(\xi\eta\zeta\,u(\xi,\eta,\zeta)=1\). The source additionally asks whether infinitely many degrees work. The allowed degrees and the meaning of ``commutator word'' remain insufficiently specified.
\subsection{Sources}
\begin{itemize}
\item Supplied local UnsolvedMath registry record: permanent ID \(700335\), exact source code \texttt{catalogue entry}, title ``Kourovka Notebook Problem 5.54.'' The record labels the problem open; this is a supplied status, not an independently verified mathematical or literature assessment.
\item \emph{The Kourovka Notebook (Unsolved Problems in Group Theory)}, original problem 5.54, fifth issue (1976), proposed by J. Wiegold. The supplied record identifies the twenty-first issue edition (2026), \texttt{arXiv:1401.0300v48}, as its source. Its collection description names E. I. Khukhro and V. D. Mazurov as editors. This bibliographic attribution is supplied and has not been independently checked.
\end{itemize}
\end{document}
